% Вставляется в main.tex после leds.tex
\section{Автономные беспилотные миссии}
\subsection*{Цели раздела}
\begin{itemize}
  \item Освоить события платформы и интерфейсы автопилота.
  \item Построить неблокирующую архитектуру миссий (таймеры, конечный автомат).
  \item Реализовать маршруты по точкам и параметрические траектории.
  \item Настроить реакции на аварии и низкое напряжение, добавить индикацию этапов.
\end{itemize}
\subsection{Введение}
Автономная миссия — заранее запрограммированная последовательность действий квадрокоптера, выполняемая без участия оператора. В рамках этого раздела мы рассмотрим модель событий платформы «Геоскан Пионер», управление полётом через функции автопилота, реакцию на внешние события и построение базовых траекторий. Документ ориентирован на академический стиль и детально объясняет каждый термин, опираясь на толковый словарь и развёрнутые примеры.

\subsection{Толковый словарь}
\begin{description}
  \item[Автопилот (\texttt{ap})] программный интерфейс управления полётом: команды взлёта, посадки, полёта к точке, обновления курса.
  \item[События (\texttt{Ev.*})] сигналы от платформы: \texttt{TAKEOFF\_COMPLETE} (взлёт завершён), \texttt{POINT\_REACHED} (точка достигнута), \texttt{COPTER\_LANDED} (посадка), \texttt{SHOCK} (удар/столкновение), \texttt{LOW\_VOLTAGE2} (низкое напряжение), \texttt{ENGINES\_DISARM} (выключение двигателей), \texttt{MCE\_*} (команды режима).
  \item[Обработчик событий (\texttt{callback(event)})] функция, которую платформа вызывает при наступлении системного события.
  \item[Таймер (\texttt{Timer})] механизм неблокирующих вызовов: \texttt{Timer.new(period, fn)} — периодический, \texttt{Timer.callLater(delay, fn)} — отложенный.
  \item[Локальная система координат] плоскость X–Y и высота Z относительно текущей базы; единицы — метры.
  \item[Курс (\texttt{yaw})] угловое направление, выражаемое в радианах; конвертация градусов: \texttt{rad = deg * math.pi / 180}.
  \item[Траектория] набор точек или правило формирования положения во времени (например, окружность).
  \item[Конечный автомат] разбиение логики на состояния с переходами по событиям; состояние хранится в переменной.
  \item[RC-каналы (\texttt{Sensors.rc})] чтение тумблеров/органов управления для запуска миссии или изменения режима.
  \item[Индикаторы (Ledbar)] световая индикация состояния миссии, аварий и этапов.
\end{description}

\subsection{Событийная модель и интерфейсы управления}
\begin{itemize}
  \item Предстартовая подготовка: \texttt{ap.push(Ev.MCE\_PREFLIGHT)} — перевод автопилота в безопасный режим подготовки. Проверяются датчики, параметры и готовность; коптер не движется. Этот шаг запускают первым, чтобы корректно инициировать взлёт.
  \item Взлёт: \texttt{ap.push(Ev.MCE\_TAKEOFF)} — команда на подъём до штатной высоты \texttt{Flight\_com\_takeoffAlt}. По достижению высоты платформа генерирует событие \texttt{Ev.TAKEOFF\_COMPLETE}.
  \item Посадка: \texttt{ap.push(Ev.MCE\_LANDING)} — команда на снижение и касание. По завершении посадки генерируется \texttt{Ev.COPTER\_LANDED}.
  \item Выключение двигателей: \texttt{ap.push(Ev.ENGINES\_DISARM)} — выключение моторов (обычно после \texttt{COPTER\_LANDED} или при аварии).
  \item Полёт к точке: \texttt{ap.goToLocalPoint(x, y, z)} — установка целевой точки в локальной системе координат (метры). По достижению цели генерируется \texttt{Ev.POINT\_REACHED}.
  \item Обновление курса: \texttt{ap.updateYaw(rad)} — установка курса в радианах; для градусов применяют перевод \texttt{rad = deg * math.pi / 180}.
  \item Считывание RC: \texttt{Sensors.rc()} — получение массива каналов радиоуправления; для тумблера \texttt{SwA} обычно анализируют 8-й канал (\texttt{rc\_chans[8]}).
  \item Обработчик событий: \texttt{function callback(event)} — функция, вызываемая платформой при наступлении \texttt{Ev.*} (\texttt{TAKEOFF\_COMPLETE}, \texttt{POINT\_REACHED}, \texttt{COPTER\_LANDED}, \texttt{SHOCK}, \texttt{LOW\_VOLTAGE2}). Внутри \texttt{callback} строятся переходы автомата миссии.
  \item Таймеры: \texttt{Timer.callLater(delay, fn)} — безопасная «задержка» без блокировки потока; \texttt{Timer.new(period, fn)} — периодическое действие через фиксированный интервал (анимации, траектории).
  \item Принципы неблокирующей логики: избегать \texttt{sleep()}; строить сценарий через события и таймеры, чтобы обработчик \texttt{callback} оставался кратким.
\end{itemize}

\subsection{Светодиоды для отладки миссии}
Светодиодная индикация — простой способ визуально подтверждать прохождение этапов миссии: изменение цвета служит отметкой «состояние автомата изменилось». Это помогает отладке без логов: по текущему цвету можно понять, выполнены ли \emph{предстарт}, \emph{взлёт}, \emph{полет}, \emph{посадка} или \emph{аварийная реакция}. В учебных примерах далее сосредотачиваемся только на полёте; использование линейки оставляем как вспомогательный инструмент для реальных испытаний.

\subsection{Справочник функций полёта}
 \begin{itemize}
   \item \texttt{ap.push(Ev.MCE\_PREFLIGHT)} — вводит борт в режим безопасной предстартовой подготовки. Проверяются датчики и параметры, движение не выполняется. Рекомендуется всегда инициировать этот шаг первым.
   \item \texttt{ap.push(Ev.MCE\_TAKEOFF)} — команда взлёта до штатной высоты \texttt{Flight\_com\_takeoffAlt}. По достижении высоты платформа генерирует событие \texttt{Ev.TAKEOFF\_COMPLETE}.
   \item \texttt{ap.push(Ev.MCE\_LANDING)} — команда на снижение и касание. Завершение посадки сопровождается событием \texttt{Ev.COPTER\_LANDED}.
   \item \texttt{ap.push(Ev.ENGINES\_DISARM)} — выключение двигателей. Применяется после \texttt{COPTER\_LANDED} или при аварийных ситуациях.
   \item \texttt{ap.goToLocalPoint(x, y, z)} — установка целевой точки в локальной системе координат (метры); оси: X — вперёд/назад, Y — влево/вправо, Z — высота. По достижении цели генерируется \texttt{Ev.POINT\_REACHED}.
   \item \texttt{ap.updateYaw(rad)} — установка курса в радианах. Для значений в градусах используйте преобразование \texttt{rad = deg * math.pi / 180}.
   \item \texttt{Timer.callLater(delay, fn)} — отложенный неблокирующий вызов функции \texttt{fn} через \texttt{delay} секунд. Удобен для последовательностей «после события — действие».
   \item \texttt{Timer.new(period, fn)} — периодический таймер; функция \texttt{fn} вызывается каждые \texttt{period} секунд, что позволяет организовать пошаговое обновление траектории/индикации.
   \item \texttt{function callback(event)} — центральный обработчик событий платформы. Внутри \texttt{callback} реализуются переходы автомата миссии и вызываются соответствующие действия.
 \end{itemize}

\subsection{Учебные ступени полёта}
\textbf{Предисловие.} Итак, мы готовы приступить к программированию квадрокоптера. В этом разделе последовательно пройдём ключевые этапы: от безопасной предстартовой подготовки до взлёта, полёта по заданным точкам и посадки. Каждый пример предваряется кратким пояснением, а строки кода снабжены комментариями, чтобы было понятно, что делает каждая команда.
\paragraph{Только запуск двигателей}
\begin{codefile}{missions/missions_engines_arm.lua}
\end{codefile}
\textbf{Пояснение.} На практике запуск двигателей выполняется в составе взлёта (\texttt{Ev.MCE\_TAKEOFF}). Явная команда \texttt{Ev.ENGINES\_ARM} полезна для проверки работы моторов без взлёта.

\paragraph{Запуск, взлёт, посадка}
\begin{codefile}{missions/missions_takeoff_land.lua}
\end{codefile}
\textbf{Пояснение.} В обработчике \texttt{callback} реагируем на завершение взлёта и инициируем посадку. После касания — выключаем двигатели.

\paragraph{Запуск, взлёт, полёт по прямой, посадка}
\begin{codefile}{missions/missions_straight_flight_simple.lua}
\end{codefile}
\textbf{Пояснение.} После взлёта выполняем полёт к точке \texttt{(x=1.0,y=0.0,z=0.8)}. По событию \texttt{POINT\_REACHED} — посадка.

\paragraph{Полет по фигуре (набор точек)}
\begin{codefile}{missions/missions_points_route.lua}
\end{codefile}
\textbf{Пояснение.} Идиоматичный «пошаговый» обход точек: каждое \texttt{POINT\_REACHED} инициирует переход к следующей цели.

\paragraph{Полет по треугольнику}
\begin{codefile}{missions/missions_triangle.lua}
\end{codefile}
\textbf{Пояснение.} Вершины равностороннего треугольника расположены на окружности: углы 0°, 120°, 240°. Переводим градусы в радианы и вычисляем координаты по \texttt{x=r*cos(a), y=r*sin(a)}.

\paragraph{Полет по квадрату}
\begin{codefile}{missions/missions_square.lua}
\end{codefile}
\textbf{Пояснение.} Квадрат центрирован в начале координат. Его вершины расположены в точках
\[
  \left(-\frac{a}{2}, -\frac{a}{2}, z\right),\ \left(\frac{a}{2}, -\frac{a}{2}, z\right),\ \left(\frac{a}{2}, \frac{a}{2}, z\right),\ \left(-\frac{a}{2}, \frac{a}{2}, z\right),
\]
что соответствует стороне длины \(\,a\,\) и фиксированной высоте \(z\). Центр квадрата — в начале координат \((0,0,z)\).

\paragraph{Полёт по окружности с поворотом квадрокоптера}
\begin{codefile}{missions/missions_circle_yaw.lua}
\end{codefile}
\textbf{Пояснение.} Два таймера управляют курсом и точкой окружности. Только неблокирующие обновления; завершение — по отложенному действию.

\subsection{Геометрия траекторий: окружность}
Окружность радиуса \(r\) в плоскости \(X\)–\(Y\) задаётся параметрически:
\[
  x(\theta) = r \cos\theta,\quad y(\theta) = r \sin\theta,\quad z = \text{const}.
\]
Здесь \(\theta\) — параметр угла в радианах. Если исходные значения заданы в градусах \(\theta_{\mathrm{deg}}\), используйте перевод в радианы
\[
  \theta_{\mathrm{rad}} = \theta_{\mathrm{deg}} \cdot \frac{\pi}{180}.
\]
Высота \(z\) выбирается постоянной для горизонтального полёта; курс \texttt{yaw} может сопровождаться как функция \(\theta\) через \texttt{ap.updateYaw}.

\subsection{Архитектурные подходы}
\begin{itemize}
  \item Конечный автомат: состояния подготовки, полёта и посадки; действия оформлены как таблица функций; переходы выполняются по событиям в \texttt{callback}.
  \item Асинхронные таймеры: пошаговое обновление логики без блокировок; одноразовые и периодические вызовы через \texttt{Timer.callLater} и \texttt{Timer.new}.
  \item Избегание блокировок: не использовать \texttt{sleep()} в сценариях полёта; для задержек применять таймеры и переходы по событиям.
  \item Безопасность: по \texttt{Ev.SHOCK} немедленно остановить таймеры и двигатели, включить аварийную индикацию; по \texttt{Ev.LOW\_VOLTAGE2} выполнить безопасную посадку или перейти в аварийный режим.
\end{itemize}

\subsection{Разбор примеров миссий}
\paragraph{Урок 12: окружность по точкам}
\begin{codefile}{missions/missions_circle_points.lua}
\end{codefile}
\textbf{Комментарий.} Генерация 12 точек окружности, переходы по событиям. Рекомендуется заменить \texttt{while/break} на явный шаг через таймер или по событию \texttt{POINT\_REACHED}. Команда \texttt{ap.push(ENGINES\_DISARM)} должна использовать пространство имён событий: \texttt{ap.push(Ev.ENGINES\_DISARM)}.

\paragraph{Урок 14: непрерывная окружность и курс}
\begin{codefile}{missions/missions_continuous_circle.lua}
\end{codefile}
\textbf{Комментарий.} Реализована непрерывная подача команды на полёт по окружности и обновление курса. Следует заменить \texttt{sleep(6)} на безопасную последовательность таймеров или переходы по событиям, чтобы не блокировать обработку задач.

\paragraph{RC-инициация миссии}
\begin{codefile}{missions/missions_rc_initiation.lua}
\end{codefile}
\textbf{Комментарий.} Запуск миссии по положению тумблера, отложенные переходы по событиям. Рекомендуется добавлять аварийные реакции на \texttt{Ev.SHOCK} и \texttt{Ev.LOW\_VOLTAGE2}.

\paragraph{Локальные точки и индикация}
\begin{codefile}{missions/missions_local_points_leds.lua}
\end{codefile}
\textbf{Комментарий.} Для корректной индикации используйте буфер размером 4, полный проход по индексам 0..3.

\subsection{Рекомендации по разработке миссий}
\begin{itemize}
  \item Делайте таймеры и состояния глобальными; избегайте локальных объектов в краткоживущих функциях.
  \item Исключайте блокировки (\texttt{sleep()}); опирайтесь на \texttt{Timer} и событийную обработку \texttt{callback}.
  \item Чётко определяйте состояния автомата; держите обработчики событий короткими.
  \item Реализуйте аварийные реакции и визуальную индикацию этапов.
  \item Работайте с курсом \texttt{yaw} в радианах; переводите градусы по формуле \(\theta_{\mathrm{rad}}=\theta_{\mathrm{deg}}\cdot\pi/180\).
\end{itemize}
